\chapter{Commands, Resources, Help, and Typed Keys}
\label{ch:commands-resources-help}
\chapterquote{A command is not a button; a resource is not a string literal; a help topic is not a tooltip accident.}{The identity rule}

\begin{chaptermap}
Core question & How should repeated screen facts keep stable identity across buttons, menus, shortcuts, labels, validation, help, resources, generated code, and tests?\\
Primary APIs & \api{ActionKey}, \api{ResourceKey}, \api{FieldKey}, \api{TextFieldKey}, \api{ComponentKey}, \api{HelpTopic}, \api{ActionDescriptor}, \api{Command}, \api{MutableCommand}, \api{CommandSet}, \api{Resources}, \api{HelpService}, \api{TooltipContent}, and \api{TooltipPolicy}.\\
Plain Swing baseline & \api{Action} already centralizes button/menu/key behavior, but applications often still duplicate strings, shortcuts, listeners, and enabled logic.\\
Swing connection & \api{SwingActionAdapter}, \api{CommandBehaviors}, key bindings, menus, toolbars, toggle buttons, F1 help, component keys, accessibility names, and composed tooltips.\\
Companion examples & \typedKeyExample, \resourcesExample, \helpServiceExample, \commandExample, \generatedActionExample, and \generationExample.\\
\end{chaptermap}

\section{Identity before presentation}

A desktop screen repeats facts. The Save operation appears in a button, a menu
item, a toolbar item, a popup menu, a key binding, a command palette, a test,
and perhaps a help page. The customer name field appears in a label, a text
field, validation problems, parse errors, help, generated tests, resources,
accessibility metadata, and screenshot fixtures. The Status column appears in a
table header, tooltip, cell renderer, validation target, column persistence
entry, export descriptor, and documentation. A screen that treats all of these
facts as unrelated strings will eventually drift.

\termdef{Stable identity}{in this chapter} is the durable name of a concept
inside the application. It is not the visible English text and not the current
component instance. The identity of Save should survive localization,
copyediting, toolbar rearrangement, menu restructuring, and a Look-and-Feel
change. The identity of a field should survive moving the field from one panel
row to another. The identity of a table column should survive a translated
header. The identity of a help topic should survive moving help content from an
HTML file to an embedded help viewer.

\termdef{Presentation}{in this chapter} is the current way an identity is made
visible or usable. A button is presentation of a command. A label is
presentation of a field. A tooltip is presentation of help, disabled state, and
problems. A keyboard shortcut is presentation of an operation for keyboard
users. A screenshot is presentation of a particular state. The presentation may
change often. The identity should change rarely and deliberately.

The distinction sounds obvious until one looks at real Swing code. A button is
constructed with \api{"Save"}. A menu item is constructed with another
\api{"Save"}. A key binding calls \api{saveCustomer}. A test finds a button by
text. A help page is opened by the string \api{"customer/save"}. The status
bar says \api{"Cannot save"}. The table column persists width under the
localized header \api{"Status"}. None of these lines is large enough to look
dangerous, but together they hide one fact behind many representations.

\termdef{Typed key}{in Corusco} is a small value whose Java type says what kind
of identity it represents. An \api{ActionKey} names an operation. A
\api{ResourceKey<String>} names a typed resource value. A \api{FieldKey}
names a semantic field. A \api{ComponentKey} names a view component contract. A
\api{HelpTopic} names help content. Most of these keys are backed by strings,
because strings are useful at resource, diagnostic, persistence, and generated
code boundaries. The type prevents the strings from being mixed accidentally.

\begin{lstlisting}[caption={Different identities, different key types.}]
static final ActionKey SAVE =
        ActionKey.of("customer/save");

static final ResourceKey<String> SAVE_TEXT =
        ResourceKey.of("customer/save/text", String.class);

static final HelpTopic CUSTOMER_FORM =
        HelpTopic.of("customer/form");
\end{lstlisting}

The stable ids remain visible. They can be searched, logged, tested, and used
at bundle boundaries. The Java type tells the compiler which family of fact the
id belongs to. Passing a help topic where a resource key is expected should not
compile. Passing a component key where a field key is expected should not
compile. This is small type discipline, but small type discipline pays off
when the same screen facts appear in generated code, runtime bindings, tests,
and documentation.

\begin{principlebox}{stable id, changing presentation}
The identity of a command, field, column, component, resource, or help topic
should survive localization, Look-and-Feel changes, and layout rearrangement.
Visible text is presentation; keys are identity.
\end{principlebox}

\begin{figure}[htbp]
\centering
\begin{minipage}[t]{0.48\linewidth}
\includegraphics[width=\linewidth]{\actionsSurfaceFigure}
\end{minipage}\hfill
\begin{minipage}[t]{0.48\linewidth}
\includegraphics[width=\linewidth]{\accessibilityLabelsFigure}
\end{minipage}
\caption{Stable identity appears through many visual surfaces. Command presenters, labels, tooltips, accessible names, and tests should refer back to typed keys rather than visible strings.}
\label{fig:commands-resources-screens}
\end{figure}

Figure~\ref{fig:commands-resources-screens} shows the ordinary visual result:
buttons, labels, tooltips, and accessible metadata. The important part is not
that the user sees keys. The user should not. The important part is that the
program has a semantic counterpart for each repeated fact. Accessibility names
should not be invented independently from field descriptors. Tooltips should
not overwrite validation state. Tests should not search for localized button
text when a command key exists. Help dispatch should not guess from component
names. Identity becomes useful precisely because it is behind the presentation.

Typed identity is also the bridge between handwritten Swing and generated
companions. Generated form metadata can expose field keys and resource keys.
Generated table metadata can expose column keys, headers, tooltips, help
topics, and persistence ids. Generated action metadata can expose action keys,
resource keys, descriptors, and descriptor lists. Handwritten code can use the
same key types when generation is not involved. The screen is then assembled
from stable facts rather than from reflection or naming conventions.

It is important not to confuse stable ids with public business identifiers.
\api{"customer/name"} is a UI or model metadata id, not necessarily a database
column, API field, or external contract. Some projects choose to align those
names, but Corusco does not require it. The stable id is stable inside the
application boundary that owns it. It should be changed deliberately because
tests, resources, generated code, and persisted UI state may refer to it.

The practical review question is simple: if this visible thing changes text,
position, or component class, what still identifies it? If the answer is
``nothing'' or ``the English label'', the screen has a hidden coupling. If the
answer is an \api{ActionKey}, \api{FieldKey}, \api{ComponentKey},
\api{ResourceKey}, \api{HelpTopic}, or column key, the screen has a durable
handle for tests, resources, generated companions, and behavior.

\section{Typed key families}

\api{ActionKey} is the stable identity for an operation. Generated
\api{@UiAction} methods create action key constants in owner-specific
companions, and handwritten code can create equivalent keys with
\api{ActionKey.of}. Action keys are used by command sets, command adapters,
menus, toolbars, keyboard bindings, diagnostics, and tests. They should name
the operation, not the visual control. A Save button and Save menu item share
one action key because they invoke one operation.

\api{ResourceKey<T>} is the typed identity for a resource value. The key
contains a stable id and a value type. In current generated UI metadata,
\api{ResourceKey<String>} is common for labels, headers, command text,
tooltips, dialog titles, and accessibility descriptions. The value type is not
decorative. A present resource with the wrong runtime type is a configuration
failure, not a value to coerce later in Swing. Resource keys let generated code
and tests agree on what resource should exist without embedding visible text in
Java source.

\api{FieldKey<O,T>} is the identity for a field belonging to an owner model.
It carries the stable field id, owner type, and value type. Text-editable field
kinds use \api{TextFieldKey}, which can be converted to a general field key
when a problem, validator, or focus resolver does not care about text-specific
editing. Field keys appear in problem targets, validation rules, generated
form companions, field resource keys, and tests. The field key is semantic; it
does not say which component currently edits the field.

\api{ComponentKey<C>} is the identity for a view component contract. It is
generic and Swing-free at the core level, so Swing code can use
\api{ComponentKey<JTextField>} or \api{ComponentKey<JTable>} without pulling
\api{java.desktop} into core. Component keys are useful when tests need to
find a component, when a problem focus resolver maps a problem to a component,
or when help behavior should be associated with a view object. Unlike a field
key, a component key is presentation-facing: it names the component contract,
not the semantic field.

\api{HelpTopic} is the stable identity for user assistance. A help topic is not
the tooltip text. A tooltip may be brief, local, and composed from several
sources. A help topic may open documentation, a support workflow, an embedded
help pane, or a browser. Generated form-field and table-column descriptors can
carry help topics from annotations. Handwritten code can use
\api{HelpTopic.of}. The id should be stable enough for help routing and
analytics.

The key families overlap in real screens. A name text field may have a
\api{TextFieldKey<CustomerEdit,String>} for validation and parsing, a
\api{ResourceKey<String>} for its label, another resource key for its tooltip,
a \api{HelpTopic} for longer documentation, and a
\api{ComponentKey<JTextField>} for tests and focus. These are not five
different opinions about the same string. They are five roles attached to one
field. Mixing them would lose information.

\begin{lstlisting}[caption={A generated-style key family for one field.}]
static final TextFieldKey<CustomerEdit, String> NAME =
        TextFieldKey.of("customer/name", CustomerEdit.class, String.class);
static final ResourceKey<String> NAME_LABEL =
        ResourceKey.of("customer/name/label", String.class);
static final HelpTopic NAME_HELP =
        HelpTopic.of("customer/name");
static final ComponentKey<JTextField> NAME_COMPONENT =
        ComponentKey.of("customer/name-field", JTextField.class);
\end{lstlisting}

The ids need not be identical, but they should be systematic. A systematic id
scheme is a maintenance tool. It lets a developer infer that
\api{customer/name/label} belongs near \api{customer/name}. It lets generated
code create resource keys predictably. It lets resource review find unused or
missing entries. It lets documentation connect field help to field metadata.
The scheme should be stable, readable, and boring.

Key ownership should be explicit. A form companion may own field keys and
field-resource keys. A presenter companion may own action keys and action
descriptors. A table companion may own column ids and table-resource keys. A
view may own component keys because components are presentation contracts. A
feature module may own help topics when help pages are organized by feature
rather than by component. The owner is not necessarily the class that first
uses the key. The owner is the place a reviewer should open when deciding
whether the key may be renamed, removed, or split. This simple convention
prevents a common generated-code problem: constants appear everywhere, but no
one knows which constant is the authoritative one.

Granularity is the next design choice. Too few keys make unrelated facts share
identity. Too many keys make the screen noisy. A useful rule is to create a
separate key when a fact can vary independently. The label of a field, the
tooltip of a field, the help topic of a field, and the component used to edit
the field can vary independently, so they deserve separate identities. The
same Save operation shown in a button and a menu item does not vary
independently, so it should use one action key. A table column and an export
column may start with the same business meaning, but if export ordering,
formatting, or visibility will be configured separately, they may need
separate ids even when they share a human label.

Avoid keys that describe an implementation accident. \api{leftPanelButton3} is
not a command. \api{mainDialogTextField2} is not a field. \api{statusColumn}
may be acceptable inside a single table companion, but \api{customer/status}
is better if the key may be used by validation, export, or help. Good keys
name the domain role at the level where the UI contract lives. They are short
enough to read in logs, long enough to avoid collisions, and plain enough that
resource authors can understand them without reading Swing code.

There is also a cost to pretending that every string is an id. A visible label
is optimized for users. It may contain punctuation, abbreviations, spaces,
case, or local grammar. An id is optimized for machines and maintainers. It
should be stable, ASCII-friendly, searchable, and insensitive to copyediting.
The two purposes conflict often enough that combining them is not economy; it
is deferred maintenance.

\begin{pitfallbox}{string-compatible is not stringly typed}
Typed keys are backed by strings so they can cross resource, persistence, and
diagnostic boundaries. That does not mean application code should pass raw
strings everywhere. Cross the boundary once, then keep the typed key.
\end{pitfallbox}

The \typedKeyExample{} companion program prints the diagnostic shape of the
baseline key families. It is intentionally simple: readers see that text field
keys, field keys, resource keys, action keys, help topics, component keys, and
problem targets all carry stable ids but remain different Java types. This is
the conceptual foundation for the rest of the chapter.

\section{Commands as operation models}

Swing's \api{Action} is one of the better parts of Swing architecture. It lets
a button, menu item, toolbar item, and key binding share one operation. It also
stores presentation properties such as name, tooltip, accelerator, mnemonic,
icon, selected state, and enabled state. A disciplined plain Swing application
can go far by using \api{Action} well.

Many Swing applications nevertheless bypass it. A button gets one listener. A
menu item gets another listener. A toolbar button gets copied text. A key
binding calls a helper. Enabled state is recalculated in each location. The
screen looks consistent on day one and drifts thereafter.

\begin{lstlisting}[caption={The duplicated-listener smell.}]
saveButton.addActionListener(event -> save());
saveMenuItem.addActionListener(event -> save());
toolbarSaveButton.addActionListener(event -> save());

saveButton.setEnabled(canSave());
saveMenuItem.setEnabled(canSave());
toolbarSaveButton.setEnabled(canSave());
\end{lstlisting}

The operation is singular, but the code does not say so. When validation logic
changes, someone must update three enabled-state branches. When the shortcut
changes, someone must remember the menu and root pane. When a test wants to
invoke Save, it may have to find a visible button. The missing object is the
operation model.

\termdef{Command}{in Corusco} is a toolkit-neutral action model. It owns a
stable \api{ActionKey}, an \api{ActionDescriptor}, observable enabled state,
observable selected state for toggle commands, and an execution method.

\begin{lstlisting}[caption={The command contract in miniature.}]
public interface Command {
    ActionKey key();
    ActionDescriptor descriptor();
    ReadableValue<Boolean> enabled();
    ReadableValue<Boolean> selected();
    boolean selectable();
    void execute();
}
\end{lstlisting}

Execution is synchronous and occurs on the calling thread. The core command
does not start background work by itself and does not marshal to Swing. When a
command is triggered by Swing, the handler normally starts on the EDT and
should delegate slow work to a task service. This rule keeps command identity
separate from task policy. The command may start a task; it should not perform
slow I/O merely because it is a command.

\begin{pitfallbox}{commands are not background workers}
A command may start a background task. It should not perform slow I/O directly
on the EDT merely because it is a command.
\end{pitfallbox}

\api{ActionDescriptor} is immutable resource-key based metadata for a command.
It stores the action key, optional text resource key, optional tooltip resource
key, optional accelerator descriptor, optional mnemonic key code, and whether
the command is selectable. This is the metadata that several UI surfaces can
share. It is not the current enabled state and not the handler.

\begin{lstlisting}[caption={Command metadata without visible strings.}]
ActionDescriptor saveDescriptor =
        ActionDescriptor.action(CustomerActions.SAVE, CustomerResources.SAVE_TEXT)
                .withTooltip(CustomerResources.SAVE_TOOLTIP)
                .withAccelerator(AcceleratorDescriptor.of(
                        KeyEvent.VK_S, InputEvent.CTRL_DOWN_MASK))
                .withMnemonic(KeyEvent.VK_S);
\end{lstlisting}

The descriptor uses resource keys rather than visible strings. A resource
resolver can later turn those keys into localized text at the Swing boundary.
Generated \api{@UiAction} methods create descriptor constants in
owner-specific companions. Handwritten presenters can build descriptors
directly. In both cases, descriptor identity is available before Swing
components exist.

\api{CommandFactory} creates mutable command instances from descriptors and
handlers. It fails fast when a selectable descriptor is passed to a non-toggle
factory or a non-selectable descriptor is passed to the toggle factory. This
failure is valuable. It prevents generated metadata and runtime behavior from
quietly disagreeing about whether selected state is part of the command.

\api{MutableCommand} is the command shape used by presenters and generated
companions when enabled or selected state changes. Save enablement may depend
on dirty state, validation problems, permissions, and busy state. A toggle
command may own selected state for a filter, preview mode, or view option. UI
controls observe the command; they should not be the authority for the command
state.

\begin{lstlisting}[caption={Enablement belongs to the command, not each button.}]
MutableCommand save = CommandFactory.command(saveDescriptor, presenter::save);

DerivedValue<Boolean> canSave = DerivedValue.of(
        () -> dirty.value()
                && !problems.value().hasErrors()
                && !busy.value(),
        dirty, problems, busy);

canSave.subscribe(event ->
        save.setEnabled(Boolean.TRUE.equals(event.newValue()), event.origin()));
\end{lstlisting}

The exact derived-value factory may vary by screen, but the ownership is the
point. The command owns enabled state. Buttons, menu items, toolbar items, and
key bindings observe that state. A test can inspect the command without
constructing a button. A command palette can list available commands without
knowing which toolbar is visible.

Toggle actions are a common source of duplication. A check-box menu item, a
toolbar toggle button, and a keyboard command may all represent the same mode.
If each control owns selected state, they drift. A selectable command exposes
selected state as an observable value. Swing adapters can copy user-originated
button selection back into the command before the handler runs. Other controls
then see the same selected state. This is the same model as ordinary command
enablement: the operation state has one owner, and presentations mirror it.

\begin{migrationbox}{one toggle owner}
When two toggle controls must stay synchronized, do not add listeners that copy
state between controls. Introduce one selectable command or value, then bind
both controls to it.
\end{migrationbox}

\api{CommandSet} groups commands under stable \api{ActionKey} values. It
preserves insertion order for deterministic menu and toolbar construction, but
lookup is by key. It owns no command lifecycle; it retains references but does
not close, reset, or disable them. Duplicate keys fail at construction time.
That failure protects tests, menus, command palettes, and generated view plans
from ambiguous operation identity.

\begin{lstlisting}[caption={A deterministic command registry.}]
CommandSet commands = CommandSet.of(save, refresh, delete);

Command refreshCommand = commands.require(CustomerActions.REFRESH);
commands.find(CustomerActions.EXPORT)
        .ifPresent(export -> menu.add(new SwingActionAdapter(export, resources)));
\end{lstlisting}

Command sets make review concrete. If a presenter exposes a command set, the
screen's operational vocabulary is visible. Hidden operations in random
listeners are easier to find because they do not appear in the set. Generated
tests can assert command state by \api{ActionKey}. Generated menu assembly can
use descriptor order. The command set is not a framework; it is a small
registry that keeps operation identity out of component fields.

Command state should be updated from model facts, not from component events
whenever possible. A Save command becomes enabled because the form is dirty,
valid, editable, and not already saving. A Delete command becomes enabled
because a durable selection exists and the current user may delete it. A
Refresh command may remain enabled during validation errors because it does
not depend on the form. These rules are not button rules. They are operation
rules. Once expressed as command state, the same facts naturally drive menu
items, toolbar buttons, shortcuts, status text, command palettes, and tests.

\api{ChangeOrigin} is useful here because command state often changes for
different reasons. User input may make Save enabled. Loading a new record may
make Save disabled because the form is clean. A background save may disable
Save while the request is running. A permission refresh may disable Delete.
The command only exposes the resulting enabled value, but surrounding code can
still propagate the origin through observable values. Tests can then
distinguish user-originated changes from programmatic refreshes when that
distinction matters.

Commands are also the right place to decide whether an operation is
contextual. Some operations depend on the current selection, but the operation
identity remains stable. \api{customer/delete-selected} may delete the current
selected row. The selected row is not part of the action key. The selection is
state read by the command handler and by the enablement rule. Other operations
are genuinely different. \api{customer/delete-selected} and
\api{customer/delete-all-filtered} should not be one command with a hidden
mode flag merely because they both use the word Delete. If a command would
need a tooltip paragraph to explain what it currently means, it may be two
commands.

Error reporting should not be hidden in the Swing action either. A command may
call a presenter method that returns a result, throws a domain exception, or
starts a task whose result is delivered later. The command adapter should not
translate all failures into modal dialogs by itself. The presenter or task
policy knows whether the result should appear as a form problem, status-line
message, dialog, retry prompt, or logged diagnostic. Keeping the adapter thin
prevents every button from inventing its own error policy.

Undo and redo fit the same pattern. The Undo command has stable identity,
enabled state, text resources, shortcut metadata, and an execution method.
The current undo presentation may need richer text, such as ``Undo Rename'',
but the operation remains Undo. A resource policy or command-description
value can contribute the changing phrase. The base action key should not
become \api{undo-rename}, \api{undo-delete}, and \api{undo-edit} unless those
are separate operations in the application command model. This distinction is
particularly important in editors, GIS tools, and design applications where
command vocabulary is large and shortcuts are learned by habit.

\section{Swing surfaces and command adapters}

\api{SwingActionAdapter} adapts one Corusco \api{Command} to a Swing
\api{Action}. It resolves descriptor metadata once during construction,
observes enabled and selected values until closed, updates Swing action
properties on the EDT, copies selected state from selectable button sources for
mutable commands, and executes the command when Swing invokes the action. The
adapter is disposable. Closing it removes subscriptions to command state. It
does not destroy, disable, or reset the command.

\begin{lstlisting}[caption={One command adapted to several Swing surfaces.}]
Command save = presenter.commands().require(CustomerActions.SAVE);
Action swingSave = new SwingActionAdapter(save, commandResources);

saveButton.setAction(swingSave);
saveMenuItem.setAction(swingSave);
toolbar.add(swingSave);
rootPane.getActionMap().put("customer.save", swingSave);
\end{lstlisting}

The adapter resolves metadata through \api{CommandResources}, a deliberately
small Swing-boundary interface for string resource keys. The adapter does not
observe later resource changes. This is a reasonable default: labels and
tooltips are normally resolved when the view is built. Applications that
support live locale switching can close and reinstall adapters or provide a
more dynamic binding. Command enabled and selected state, however, is observed
while the adapter is installed.

Do not share one adapter instance indiscriminately across components with
different lifecycles. Swing permits one \api{Action} to be used by several
components, and that is often convenient inside one view. But the adapter owns
subscriptions to command state. If a toolbar outlives a dialog button, the
shared adapter's lifecycle becomes unclear. A generated behavior plan can
install one adapter per component and close it with the component's behavior
scope. The command remains shared; the Swing adapter is a view binding.

\api{CommandBehaviors} wraps this pattern for generated and handwritten
behavior plans. \api{commandButton} installs a command action on a button-like
component and restores the previous action on close. \api{commandMenuItem}
names the same pattern for menus. \api{commandKeyBinding} installs a command's
accelerator into an input map and action map, records previous entries, and
restores them on close. These behaviors are lifecycle handles, not global
registration magic.

\begin{lstlisting}[caption={Command behavior installation is reversible.}]
try (BehaviorScope scope = new BehaviorScope()) {
    scope.install(saveButton,
            List.of(CommandBehaviors.commandButton(save, resources)));
    scope.install(rootPanel,
            List.of(CommandBehaviors.commandKeyBinding(
                    save, resources,
                    JComponent.WHEN_ANCESTOR_OF_FOCUSED_COMPONENT)));
}
\end{lstlisting}

The key binding uses the command descriptor's accelerator metadata. The
operation identity remains the \api{ActionKey}. The Swing input-map key may be
the \api{ActionKey} itself or a Swing-local value; it is an implementation
detail of that view. The important rule is that keyboard invocation and button
invocation call the same command and observe the same enabled state. A shortcut
that works while the button is disabled is a symptom of duplicate operation
paths.

\begin{figure}[htbp]
\centering
\begin{minipage}[t]{0.48\linewidth}
\includegraphics[width=\linewidth]{\actionsFocusFigure}
\end{minipage}\hfill
\begin{minipage}[t]{0.48\linewidth}
\includegraphics[width=\linewidth]{\behaviorPlanFigure}
\end{minipage}
\caption{Command identity participates in keyboard workflow and generated behavior plans. The same command metadata should drive visible controls, shortcuts, and installed interaction behavior.}
\label{fig:commands-behavior-plan}
\end{figure}

Figure~\ref{fig:commands-behavior-plan} connects commands to the surrounding
view plan. Keyboard shortcuts are not loose conveniences. They are part of the
user workflow and must share command state. A generated behavior plan is not a
bag of listeners. It is an ordered installation of reversible behaviors:
command binding, validation decoration, composed tooltip, F1 help, busy
overlay, and accessibility metadata. Command identity is one participant in
that plan.

Disabled reasons deserve explicit modeling. Enabled state tells whether a
command can run. A disabled reason explains why. Swing \api{Action} has no
standard disabled-reason property. Many applications solve this by overwriting
the tooltip whenever the command becomes disabled. That works until static
help, validation problems, and help hints need the same tooltip. Model the
reason as data, such as a \api{ReadableValue<String>} or a problem set. Then a
tooltip policy, status line, command palette, or accessibility description can
compose it with other contributors.

\begin{lstlisting}[caption={Command enablement with a reason value.}]
SimpleValue<String> saveDisabledReason = SimpleValue.of("");

canSave.subscribe(event -> {
    boolean enabled = Boolean.TRUE.equals(event.newValue());
    save.setEnabled(enabled, event.origin());
    saveDisabledReason.setValue(enabled ? "" : "Fix validation errors",
            event.origin());
});
\end{lstlisting}

The reason is not merely visual text. It is part of the command state from the
user's point of view. Save disabled because there are validation errors is
different from Save disabled because there are no changes, because the record
is read-only, because a save is already running, or because permissions are
missing. A boolean cannot express those differences. A reason value can.

Swing's input maps make duplicate paths especially easy to create. A root pane
can bind Ctrl+S to one action, a focused editor can bind Ctrl+S to another
action, a menu item can carry an accelerator, and a toolbar button can listen
directly to a helper. Swing will dispatch according to focus and input-map
condition, not according to the developer's intent. Corusco does not change
Swing's dispatch rules. It makes the installed target explicit. A command key
binding should be installed with the intended input-map condition, and the
action-map entry should invoke the same command object as the visible
presentations.

The input-map condition is a design decision. \api{WHEN\_FOCUSED} belongs to
component-local editing gestures. \api{WHEN\_ANCESTOR\_OF\_FOCUSED\_COMPONENT}
belongs to a panel or dialog workflow. \api{WHEN\_IN\_FOCUSED\_WINDOW} belongs
to broader window commands and should be used carefully because it can shadow
local editing behavior. A generated behavior plan can record the condition
next to the descriptor. A code review can then ask whether the shortcut should
belong to the field, the form, the dialog, or the whole window. Without that
explicit question, shortcuts tend to migrate upward until they become
surprising.

Menus add another wrinkle. A \api{JMenuItem} created from a Swing
\api{Action} reads standard action properties, but menus also participate in
selection models, popup lifecycles, and platform menu conventions. A popup
menu attached to a table row may need to refresh command enablement just
before it opens because selection changed under the pointer. That refresh
should update command state, not replace the menu item's listener. The item
then remains a presentation of the command, and the popup merely asks the
presenter to re-evaluate context before showing.

Toolbars have density pressure. A toolbar may omit text, show only icons, or
group commands visually. This is exactly where command descriptors help. The
toolbar can choose an icon-only presentation while the command descriptor
still supplies text for accessible names and tooltips. If the toolbar stores
only an icon and a listener, accessibility and tests are forced to guess. If
it adapts a command, the visible density can change without losing semantic
metadata.

Do not confuse adapter reuse with command reuse. The command is the shared
operation model. Adapters are Swing objects installed into a particular view
surface. A dialog can create a button adapter and a key-binding adapter for
the same command, then close both when the dialog closes. A main window can
keep its own adapter for the same command if the command itself is shared at
the application level. This keeps lifecycle local without duplicating
operation state.

\begin{pitfallbox}{two operation paths}
If a shortcut can run while the button is disabled, or if a menu item and
toolbar button disagree about selected state, the screen probably has duplicate
operation paths. Route every presentation through the same command.
\end{pitfallbox}

\section{Resources and localization boundaries}

\termdef{Resource key}{in Corusco} is a typed key for a value resolved from a
resource source. A \api{ResourceKey<String>} may name a label, tooltip, table
header, dialog title, command text, accessibility description, or help hint.
The \api{Resources} interface resolves keys and checks the expected value type.
It is Swing-free and does not prescribe whether values come from a map, a
resource bundle, a database, a generated file, or an application localization
service.

\begin{lstlisting}[caption={Typed resource lookup.}]
Resources resources = Resources.of(Map.of(
        "customer/save/text", "Save",
        "customer/save/tooltip", "Save customer changes"));

String text = resources.require(CustomerResources.SAVE_TEXT);
String tooltip = resources.resolve(CustomerResources.SAVE_TOOLTIP, "");
\end{lstlisting}

The boundary is important. Inside the application, generated descriptors hold
typed resource keys. At the resource provider boundary, a map or bundle may be
keyed by strings. Lookup checks that a present value is assignable to the
key's declared value type. Missing optional values are represented as
\api{Optional.empty} or fallbacks. Required lookup throws \api{ResourceException}
when absent or wrong-typed. This contract catches configuration mistakes close
to the resource boundary rather than later in component construction.

\api{MapResources} is intentionally small. It copies an input map, treats ids
as stable resource ids, rejects null values, and checks runtime type during
lookup. It is suitable for examples, tests, and small applications. A
production application may have a localized, reloadable, cached, or layered
resource service. That service can still implement \api{Resources}; generated
descriptors and tests do not need to know the source.

Resource failure policy has two useful modes. Required lookup fails when a
resource must exist in a finished application. Fallback lookup supplies a
default when missing text is acceptable. Use required lookup aggressively in
tests for generated UI text, menus, table headers, and dialog titles. A missing
resource is cheap to fix before release and annoying to discover in an
installer. Use fallbacks for optional help, experimental examples, diagnostics,
or soft metadata where absence is part of the policy.

\begin{pitfallbox}{localized text is not identity}
Do not use a visible table header, button text, or tooltip as a persistence
key, test selector, command id, or help topic. Localization and copyediting will
break it.
\end{pitfallbox}

Generated resource companions solve a problem that old Swing applications
often solve poorly. One extreme hard-codes strings in component construction.
The other extreme uses a resource bundle whose keys are unrelated to Java
declarations. Generated companions keep identity near source declarations while
keeping visible strings external. A generated table companion may expose header
and tooltip keys for each column. A generated form companion may expose label,
tooltip, and help keys for each field. A generated action companion may expose
text and tooltip keys for each command.

\begin{coruscobox}{Resource keys from declarations}
Corusco generation makes resource identity follow the source declaration. The
generated key is stable, typed, and available to tests, while the visible
string remains external.
\end{coruscobox}

Localization is not only translation. It is also review, terminology,
accessibility wording, and product tone. A label may change because the product
team chooses clearer language. A tooltip may be shortened. A command text may
be platform-specific. A mnemonic may differ by locale. If tests, persistence,
or help routing depend on the visible text, ordinary copyediting becomes a
breaking change. Resource keys keep copyediting in the presentation layer.

Resource keys can also support consistency checks. A build or test can scan
generated resource keys and assert that required values exist. It can compare
resources across locales. It can detect duplicate visible text if that matters.
It can ensure every visible command has tooltip text or every help-enabled
field has a static hint. These checks are possible because the resources are
named facts, not incidental strings.

The \resourcesExample{} companion program demonstrates typed lookup with
generated table resource keys. It resolves a header, resolves a tooltip, and
uses a fallback for a missing optional value. That is the intended scale of the
resource abstraction: small enough for examples, strong enough to keep
generated metadata and runtime presentation aligned.

Resource bundles remain a fine storage mechanism. The point is not to replace
\api{ResourceBundle} with a new religion. The point is to keep the typed
boundary at the place where application code asks for a value. A provider may
load bundle keys, layer customer-specific overrides, cache values, and expose
them through \api{Resources}. Generated descriptors and Swing adapters still
see typed keys. Tests can still ask for required keys. The underlying storage
may change from bundles to database-backed localization without rewriting
button construction.

Pseudo-localization is a useful resource test for Swing. Replace English
strings with longer bracketed strings, accented variants, or deliberately wide
phrases, then run the screenshot harness. MigLayout forms that use sensible
growth constraints usually survive. Hand-positioned labels, narrow buttons,
and table headers that assume English length do not. This kind of test belongs
near resource infrastructure because it tests the same boundary: visible text
may change while identity stays stable. The harness should locate controls by
typed keys, not by the pseudo-localized text it is trying to inspect.

Mnemonic policy also belongs at the resource boundary more often than early
examples suggest. English examples often write \api{KeyEvent.VK\_S} beside the
Save descriptor, and that is acceptable for a small fixed-language program.
In a localized application, visible text and mnemonic choice are reviewed
together. The descriptor may carry a default mnemonic, but a locale resource
or platform policy may override it. The command key remains
\api{customer/save}. The visible text and mnemonic remain presentation.

Accelerators differ from mnemonics because users learn them as commands, not
as underlined letters in localized labels. Ctrl+S for Save is often stable
across locales. Even so, accelerator policy may vary by platform, by product
edition, by input mode, or by user preference. A descriptor is a good place
for the default accelerator. The behavior installation is the place to apply
application policy and decide where the key is active. Tests should assert the
default descriptor when generation is under test and assert installed bindings
when a view's workflow is under test.

Resource values can be richer than strings, but the first version of most UI
metadata should resist overgeneralization. Strings, key strokes, icons, colors,
and format patterns all have different review processes. A typed
\api{ResourceKey<T>} can name any value type, but mixing all presentation
assets into one unstructured map tends to hide policy. A command icon, for
example, may need light and dark variants, high-contrast variants, and
accessibility review. It may be better represented by an icon service keyed by
action identity than by a raw \api{ResourceKey<Icon>}. The type system enables
the choice; it does not remove the need for design ownership.

Resource tests should be written at several levels. A unit test for
\api{MapResources} verifies missing and wrong-typed values. A generated
metadata test verifies that every required descriptor key exists in the sample
resource set. A screen smoke test verifies that the resolved values reach
buttons, labels, table headers, and tooltips. A screenshot review verifies
that the resulting text fits. These are not redundant tests. Each one protects
a different boundary: provider correctness, metadata completeness, Swing
binding, and visual layout.

\section{Help topics and composed tooltips}

\termdef{Help topic}{in Corusco} is a stable identifier for user assistance.
It is not the tooltip text. A tooltip may be short, local, and transient. A
help topic may open a page, focus an embedded documentation panel, launch a
support workflow, record analytics, or invoke a platform help system. The
topic id is the durable bridge between screen metadata and whatever help
infrastructure the application chooses.

\api{HelpService} accepts a \api{HelpTopic} or \api{HelpRequest}. A request
can carry the source object and optional context. The default help service is
simple: it records the most recent request for diagnostics and delegates to a
configured handler. A handler might open a browser, show embedded help, or
record the request in a test. The core service owns neither the Swing
component in the source field nor the UI opened by the handler.

\begin{lstlisting}[caption={Opening typed help.}]
HelpService help = application.helpService();

help.open(CustomerHelp.CUSTOMER_FORM, formPanel, "credit-limit");
\end{lstlisting}

The common Swing help gesture is F1. But F1 should not simply open a generic
page. It should open the topic associated with the focused component, selected
command, table column, or screen. That requires metadata. A generated field
descriptor can carry a help topic. A component key can identify the focused
component. A table column descriptor can carry column help. A behavior can walk
from component to descriptor to service request.

\api{StandardBehaviors.helpOnF1} shows the Swing boundary. Installation
requires a help service in the behavior context. The behavior installs an F1
input/action-map entry on the component, invokes the help service from the EDT,
and restores previous entries during cleanup. The help behavior does not know
where help content lives. It routes a typed topic to the configured service.

\begin{lstlisting}[caption={Help behavior is metadata plus reversible Swing wiring.}]
BehaviorScope scope = new BehaviorScope(helpService);

scope.install(nameField,
        List.of(StandardBehaviors.helpOnF1(CustomerHelp.NAME_FIELD)));
\end{lstlisting}

The behavior's cleanup is part of the same lifecycle story as command key
bindings. Installing F1 help changes Swing input and action maps. Closing the
behavior restores previous entries. A help binding left behind after a dialog
closes can route F1 to a dead component or stale topic. Help is user-facing,
but the implementation is still a listener relationship with an owner.

Tooltips often have several contributors: validation problems, parse problems,
disabled reasons, static descriptor help, formatting examples, and a help
available hint. Swing components have one tooltip slot. If each contributor
calls \api{setToolTipText}, the last writer wins and the user loses context.
\api{TooltipContent} keeps the contributors separate. \api{TooltipPolicy}
composes them in deterministic order.

\begin{lstlisting}[caption={Composing tooltip text from several concerns.}]
TooltipContent content = new TooltipContent(
        form.problems().filter(ProblemFilter.field(CustomerFields.NAME)),
        disabledReason,
        resources.resolve(CustomerResources.NAME_TOOLTIP, ""),
        true);

Optional<String> tooltip = TooltipPolicy.standard().compose(content);
\end{lstlisting}

The standard policy places the most severe problem first, then the disabled
reason, then static help, then the help indicator. It returns ordered
plain-text lines or a newline-separated string. It does not perform resource
lookup, localization, problem filtering, or Swing HTML rendering. This is
intentional. A renderer, binding, accessibility adapter, or test can inspect
the same ordered content. Swing-specific formatting remains at the edge.

The ordering is humane. If a field is invalid, the first thing the user needs
is the problem. If the command is disabled, the user needs the reason. If
nothing is wrong, static help can explain the expected input. If richer help is
available, the final line can advertise F1. A tooltip that always starts with
static help hides urgent feedback. A tooltip that always overwrites help with
disabled text hides documentation. A composed tooltip lets each concern keep
its place.

Help and tooltip metadata also support accessibility. A screen reader may need
an accessible description rather than a hover tooltip. A keyboard user may need
F1 help because hover is irrelevant. A screenshot test may need stable text
for a tooltip state. Keeping problem messages, disabled reasons, static help,
and help topics separate gives the application several ways to present the
same underlying facts. Collapsing them into one string gives only one fragile
presentation.

The \helpServiceExample{} companion program demonstrates help dispatch through
generated descriptor metadata. It does not open a browser because the point is
the contract: a descriptor supplies a stable topic, the service records and
dispatches a structured request, and the test asserts topic id and context.
The visual help system can evolve after that boundary is stable.

Help context should be chosen deliberately. A field topic may need a field id,
the current validation problem, or the current parse text. A table-column
topic may need the row type and column id but not the selected row's private
data. A command topic may need the action key and disabled reason. Passing no
context is sometimes correct; passing the entire presenter is usually too
much. A \api{HelpRequest} should contain enough information for the help
handler to route the user, not enough information to become a hidden service
locator.

Table help is a good stress test. The focused component is only the
\api{JTable}, while the user's question is often about a cell, row, or column.
The help behavior therefore needs help metadata from the table descriptor and
selection or hit-test context from Swing. A generated column descriptor can
carry the help topic. The installed behavior can determine the current column
and submit a request for that topic with column context. The help service does
not need to know how \api{JTable} calculates view and model coordinates. It
receives a structured request.

Command help is similar. F1 while a toolbar button is focused may open help
for the command, not for the button class. A disabled command can route help
that explains the operation and perhaps the disabled reason. A command palette
can show a help affordance for each listed command because action identity and
help identity are explicit. Without command identity, help systems often fall
back to page-level help, which is cheaper to wire but less useful when the
user is stuck on a particular operation.

Tooltip text has a short shelf life. It is read in a moment, often while the
pointer is covering the screen. It should not become a dumping ground for all
documentation. If the tooltip must contain several lines of explanation, ask
whether the static help belongs in a help topic instead. A composed tooltip is
valuable because it can put urgent transient facts first and leave longer
material to F1 help. The tooltip policy is therefore also an editorial policy:
short problem, short reason, short hint, longer help elsewhere.

Problems and disabled reasons should not fight each other. A field can be
invalid while the Save command is disabled because of that invalid field. The
field tooltip should lead with the field problem. The Save button tooltip may
lead with ``Fix validation errors''. A problem summary may list all invalid
fields. These are different presentations of related facts. If every
presentation pulls from the same problem set and command reason values, they
can remain consistent without sharing one giant text string.

Help availability must also be discoverable to keyboard users. A hover-only
hint is insufficient. A visible Help button, F1 binding, menu item, command
palette entry, or accessible description can expose the same topic. Corusco's
help topic type does not prescribe which surfaces are present. It gives each
surface a stable target. The UI design can then decide how explicit help
should be for the screen's audience and density.

\section{Generated metadata and handwritten code}

Corusco generation is most valuable when it turns repeated facts into stable
metadata without hiding the screen's design. An \api{@UiAction} method can
produce action keys, text resource keys, descriptors, selectable flags,
mnemonics, accelerators, and descriptor lists. A \api{@SwingForm} model can
produce field keys, resource keys, validation metadata, and generated binding
helpers. A \api{@SwingTable} descriptor can produce column keys, headers,
tooltips, help topics, persistence ids, and table helpers. The generated code
does not need runtime reflection because the facts are ordinary constants and
methods.

\begin{lstlisting}[caption={Generated action metadata is ordinary Java.}]
@UiAction(
        id = "generated-customer/save",
        tooltip = "generated-customer/save/tooltip",
        mnemonic = 83)
void save() {
}

ActionDescriptor descriptor =
        GeneratedActionMetadataExampleActions.SAVE;
\end{lstlisting}

Generated metadata should stop at the right boundary. A descriptor is not a
button. A resource key is not a localized string. A help topic is not a help
viewer. A field key is not a text field. A generated companion can provide
these facts so handwritten or generated view assembly can use them. The screen
author still decides layout, grouping, density, and workflow. This is the same
theme as in the dialog and MigLayout chapters: generation should remove
repeated wiring, not bury product design.

Handwritten code should use the same types. A feature that is not generated
can still declare an \api{ActionKey}, \api{ResourceKey}, and
\api{ActionDescriptor}. A custom component can still have a
\api{ComponentKey}. A one-off help topic can still be a \api{HelpTopic}. This
keeps handwritten and generated parts interoperable. Tests do not need to know
whether a key came from an annotation processor or a handwritten companion.

\begin{coruscobox}{Generated facts, visible design}
Corusco generation should produce stable facts and safe wiring: keys,
descriptors, resource identities, validation metadata, and installable
behaviors. Layout, grouping, and workflow remain reviewable application code
unless the project explicitly chooses generated layout for a narrow purpose.
\end{coruscobox}

Generated names should be predictable. A presenter action companion might be
named after the presenter. A table resource companion should be named after the
row or table descriptor. A form fields companion should be named after the
form. Predictability lets chapters, tests, and examples mention generated
artifacts without long paths. It also lets developers navigate by intent:
field metadata near the form, action metadata near the presenter, table
metadata near the row descriptor.

Annotation processor diagnostics should use the same identity vocabulary.
Duplicate action ids, missing resource ids, unsupported mnemonic values,
invalid help topics, and conflicting generated names should be reported before
runtime. A runtime failure in a generated Swing adapter is late; a compiler
diagnostic is early. This is one reason generated metadata is preferable to
runtime annotation scanning for Swing screens: the build can fail while the
source context is still present.

Generated descriptor lists can assemble menus and toolbars deterministically,
but application policy still decides which descriptors appear where. A menu may
include every command. A toolbar may show only frequent commands. A popup menu
may show row-specific commands. A command palette may list hidden operations.
All of these surfaces can share one command set and descriptor identity while
choosing different presentations.

The \generatedActionExample{} companion program demonstrates generated action
metadata from annotated methods. The \generationExample{} companion program
shows generated customer metadata in a broader example. Read both as identity
examples before reading them as generation examples: which ids are stable,
which strings are resources, which descriptors are generated, and where
runtime command state begins.

\section{Testing identity and behavior}

Typed keys make tests less brittle. A test can assert that the Save command is
disabled, that it has the expected action key, that its text resource resolves,
that its tooltip resource exists, and that F1 opens the expected help topic.
It does not need to search for a button whose current localized text is
\api{"Save"}. Visible text can be tested where visible text matters, but it
should not be the only handle for a command.

\begin{lstlisting}[caption={Testing command state by key.}]
Command save = commands.require(CustomerActions.SAVE);

assertThat(save.key()).isEqualTo(CustomerActions.SAVE);
assertThat(save.isEnabled()).isFalse();

form.name.setRawText("Grace", ChangeOrigin.USER);
assertThat(save.isEnabled()).isTrue();
\end{lstlisting}

The Swing adapter deserves a narrow EDT test. Construct a command, adapt it,
install it on a button or menu item, change command enabled state, and assert
the Swing component updates. For selectable commands, trigger a toggle button
and assert that the mutable command selected state receives a user-originated
change. Then close the adapter or behavior scope and assert that later command
changes do not update the disposed component. This tests Swing binding
lifecycle without testing business policy through button clicks.

Resource tests should assert required resources through typed keys. A generated
form resource test can iterate generated required label keys and call
\api{resources.require}. A generated table resource test can require every
header key. A command resource test can require every command text key and
optionally every tooltip key. Wrong-typed values should be tested at least once
for the resource provider. Missing resources should fail before screenshots or
manual testing reveal blank labels.

Help tests should assert structured requests. A fake or default help service
can record the last \api{HelpRequest}. An F1 behavior test can install help on
a component, trigger the input/action-map entry, and assert that the request
contains the expected \api{HelpTopic}, source component, and context. The test
does not need a browser. It proves the boundary between Swing focus behavior
and application help routing.

Tooltip tests should inspect composed lines, not only final Swing strings. A
policy test can pass problems, a disabled reason, static help, and help
available, then assert the ordered lines. A Swing binding test can separately
assert that the component's tooltip receives the composed string and is
restored on close. This separation avoids coupling core tooltip policy to
Swing formatting.

Component-key tests should prefer typed markers over component names.
\api{SwingComponentKeys.mark} stores a \api{ComponentKey} in a client property.
\api{SwingMvpTester} can then find components by key, enter text, select
tables, execute commands by \api{ActionKey}, and assert problems by
\api{FieldKey}. Component names remain a useful compatibility fallback, but
typed keys communicate what kind of component the test expects.

\begin{figure}[htbp]
\centering
\includegraphics[width=0.86\linewidth]{\testingKeysFigure}
\caption{Typed component keys make Swing tests speak about generated view contracts instead of localized text or fragile component traversal.}
\label{fig:commands-testing-keys}
\end{figure}

Figure~\ref{fig:commands-testing-keys} belongs in this chapter even though the
testing chapter returns to it later. Testing is one of the clearest payoffs of
stable identity. A test that says \api{executeCommand(CustomerActions.SAVE)}
or \api{enterText(CustomerViewKeys.NAME\_FIELD, "Grace")} is easier to read
and less fragile than a test that searches for a button named \api{"Save"} and
a text field at index two.

Screenshot tests need both stable identity and visible text. The screenshot
should show localized labels and real tooltips. The test harness should locate
components and command states by typed keys. If a screenshot changes because
the label text was improved, the screenshot diff should show the visual
change, but the harness should not fail to find the component. Stable identity
keeps visual review and test mechanics separate.

\section{Worked review of a customer screen}

Consider a customer screen with a search field, a results table, a details
form, Save and Refresh commands, a Delete command available only for selected
rows, a Show inactive toggle, field-level help, and generated table columns.
The plain Swing implementation may start with button text, menu text, and
listeners. The Corusco review starts by listing identities.

The operations are \api{customer/save}, \api{customer/refresh},
\api{customer/delete}, and \api{customer/show-inactive}. Each receives an
\api{ActionKey}. Save, Refresh, and Delete are ordinary commands. Show inactive
is selectable. Each has a descriptor with resource keys for text and tooltip.
Refresh may have an accelerator. Save may have a mnemonic. Delete may have a
disabled reason when no row is selected or permissions are missing.

The fields are \api{customer/name}, \api{customer/credit-limit}, and
\api{customer/status}. Each receives field keys through generated form
metadata. Labels and tooltips are resource keys. Field help topics point to
documentation. Components receive component keys for tests and focus recovery:
name field, credit field, status combo, search field, results table. The table
columns have column keys, header resource keys, tooltip resource keys, help
topics, and persistence ids. None of these ids is the English text shown to the
user.

The presenter constructs commands from descriptors and handlers. Save
enablement derives from dirty state, problem state, busy state, and permissions.
Refresh enablement derives from screen lifecycle and refresh busy state. Delete
enablement derives from selection and permissions. Show inactive selected
state is a model value. The button row, menu, toolbar, popup menu, and key
bindings all adapt the commands. There is no separate delete listener on the
popup menu.

Resources are verified before the screen is considered complete. Every command
text key resolves. Required field labels resolve. Table headers resolve. Static
tooltips resolve where descriptors declare them. Optional help hints use
fallbacks only where the policy says optional. The resource test uses typed
keys and does not assert that English text is unchanged unless the text itself
is the product requirement.

Help behavior installs F1 on fields and the results table. If focus is in the
name field, F1 opens the name help topic. If focus is in the table header or a
cell whose column descriptor has help, F1 opens column help. If focus is in the
screen background, F1 opens screen help. The help service receives structured
requests with topic, source, and context. Tooltips advertise help availability
only when a topic exists.

Tooltips are composed. The name field's tooltip may show ``Name is required''
when invalid, then a disabled reason if the field is read-only, then static
help, then an F1 hint. The Save button's tooltip may show why Save is disabled
or the static save tooltip. The table column tooltip may show a column
description and help indicator. No part of the screen should have three
listeners racing to call \api{setToolTipText}.

Tests speak in identities. A presenter test executes Save by \api{ActionKey}
and asserts that the handler did not run while disabled. A resource test
requires generated keys. A help behavior test triggers F1 and asserts the
topic. A Swing view test finds the name field by \api{ComponentKey}. A table
state test persists column width under a stable column id, not the localized
header. The tests become a map of the screen's semantic structure.

The review table is short but effective:

\begin{center}
\textbf{Command and identity review table.}
\end{center}

\noindent\begin{tabularx}{\linewidth}{@{}p{0.32\linewidth}X@{}}
\toprule
Symptom & Review question\\
\midrule
Two listeners call the same method & Should this be one command presented twice?\\
Button and menu enabled state differ & Which object owns command enablement?\\
Shortcut works while button is disabled & Is the key binding using the same command?\\
Tooltip disappears when disabled & Which contributors should the tooltip compose?\\
Test searches for button text & Can the test locate the command by \api{ActionKey}?\\
Localized table header used as id & Which typed key or column key is stable?\\
F1 opens a generic page & Which descriptor supplies the focused help topic?\\
Generated metadata is hard to find & Are companion names predictable and close to declarations?\\
\bottomrule
\end{tabularx}

The table is not meant to shame plain Swing. Plain Swing \api{Action} is good
architecture. The Corusco question is what happens around it: typed identity,
resource lookup, generated descriptors, help routing, composed tooltips,
lifecycle cleanup, and tests. Corusco supplies names for those surrounding
concerns so they do not become conventions carried only in memory.

\section{Operational conventions}

Identity systems fail when conventions are too clever to remember. A project
should choose boring id shapes and keep them. Action ids might use
\api{area/object/verb}, such as \api{customer/save} or
\api{orders/export}. Field ids might use \api{object/field}, such as
\api{customer/name}. Resource ids may extend the semantic id with a role:
\api{customer/name/label}, \api{customer/name/tooltip},
\api{customer/save/text}. Help topics may share the semantic stem when that
helps documentation: \api{customer/name}. Component keys may include a view
role: \api{customer/name-field} or \api{customer/results-table}. The exact
scheme is less important than consistency, but consistency should be checked
by code review and generation tests.

Names should be stable, not overly encoded. Avoid ids that include visual
position, English words, row numbers, or Swing class names unless those facts
are truly part of the identity. \api{customer/edit-dialog/row-3/name} is a
bad field id because the row may move. \api{customer/name} is better.
\api{customer/nameTextField} may be a reasonable component key in a small
handwritten view, but \api{customer/name-field} is less tied to the current
component class. \api{saveButtonText} is a poor resource id because the Save
operation may appear in menus and command palettes too.

The id scheme should also leave room for sibling concepts. A field may need a
label, tooltip, placeholder, accessible description, validation message, and
help topic. If the first resource key is \api{customer/name}, there is no
obvious place for siblings. If the field id is \api{customer/name} and the
label resource is \api{customer/name/label}, the family can grow. Generated
companions should follow the same pattern so handwritten keys and generated
keys feel like one vocabulary.

Versioning is a real concern. If a key is used only inside one screen and no
state is persisted under it, changing the id is a refactoring. If a key is used
for table-column persistence, saved layouts, user preferences, screenshots,
external help routing, or analytics, changing it is a migration. The code does
not know the difference unless the project records where keys cross durable
boundaries. Table column ids and action ids often become durable sooner than
developers expect. Treat them with the same respect as serialized field names.

Resource audits should be automated. A generated companion can expose the
resource keys it requires. A test can iterate them and call
\api{resources.require}. A localization build can compare keys across locales.
A screenshot build can run with a pseudo-locale whose strings are longer, so
layout problems appear before real translations arrive. A resource audit can
also detect resource values that are present but wrong-typed. These checks are
cheap because resource identity is typed and enumerable.

Missing resource policy should be different in development, tests, and
production only when the difference is deliberate. Tests should generally fail
fast on missing required resources. Development builds may show visible marker
text for missing optional strings so the screen can still be inspected.
Production builds may fail for missing required command text but quietly omit
optional help hints. What matters is that the policy is in the resource
boundary, not scattered across button construction and table header code.

Live locale switching is possible, but it should not be accidental. The simple
\api{SwingActionAdapter} resolves metadata once during construction. That is
excellent for ordinary applications. If the application supports live resource
replacement, commands and components need resource bindings rather than
one-time resolution. The view can close and reinstall adapters, or a dynamic
resource binding can observe resource changes. Do not assume one-time
adapters will magically update when a resource bundle changes. The adapter
contract should remain explicit.

Permissions are another source of command-state confusion. A Save command may
be disabled because the form is invalid, because the form is clean, because a
save is running, or because the user lacks permission. These reasons are not
equivalent. Permission failure may need a lock icon, a tooltip, a help link, or
an audit message. Validation failure needs problem focus. Busy state may need a
progress indicator. Clean state may need no explanation. Modeling only one
boolean loses the difference. A command can expose enabled state while a
neighboring value or richer wrapper exposes the reason.

Command palettes make identity discipline visible. A command palette lists
operations independently from their current buttons. It needs action keys,
localized names, shortcuts, enabled state, selected state for toggles, and
perhaps disabled reasons. It should not scrape buttons from the current view.
A palette built from command sets and descriptors naturally includes commands
that have no toolbar item. A palette built from Swing components can only show
what is currently visible and risks invoking stale listeners.

Menus and toolbars have different density and discovery rules. A toolbar
should show frequent commands with concise labels or icons. A menu can expose
more complete command vocabulary. A popup menu should be contextual. These
differences are presentation policy. They should not create new command
objects unless the operation is genuinely different. If Delete in the toolbar
and Delete in the table popup have different enablement rules, ask whether
there are two operations or one operation with a contextual disabled reason.

Keyboard shortcuts should be reviewed as part of command metadata, not as
afterthoughts. The accelerator in a descriptor is user-facing contract. It may
vary by platform, input map condition, or application mode. A generated
descriptor can carry the default shortcut; application policy can override it
when installing behaviors. The key binding should still invoke the same
command. A shortcut installed directly to a helper method is a second command
path without command state.

Mnemonics deserve a smaller but similar review. A mnemonic belongs to a
localized presentation, not the operation identity. In English, Save may use
\api{S}. In another language, the visible text and mnemonic may differ. A
descriptor may carry a default mnemonic key code, but a resource or locale
policy may override it. Tests should not assume the English mnemonic unless
the test is specifically for the English resource set.

Icons are presentation too. The current \api{ActionDescriptor} focuses on
resource-key metadata, accelerator, mnemonic, and selectable state. An
application may add an icon policy around descriptors or through a richer
resource layer. The same principle holds: the icon should be associated with
the command identity, not copied independently into each button. If a toolbar
uses only icons, accessible name and tooltip resources become even more
important because visible text is absent.

Accessibility review should treat typed identity as source material, not final
output. A field label resource can help set the accessible name for the
component. A tooltip or static help resource can help form an accessible
description. A disabled reason can explain why a control cannot be used. A help
topic can provide a longer route. The accessibility layer may need different
phrasing from visual labels, so do not blindly reuse every visible string. But
do derive the relationship from the same keys and descriptors so accessible
metadata does not drift away from the screen.

Diagnostics should print typed identities. A log entry that says ``Save failed''
is less useful in a localized application than one that says
\api{ActionKey[customer/save]} and then includes the localized text if useful.
A validation diagnostic should include the \api{FieldKey}. A help dispatch
diagnostic should include the \api{HelpTopic}. A missing resource diagnostic
should include the resource id and expected value type. Typed keys make such
diagnostics natural because their \api{toString} output already names the
family.

Analytics should use stable ids, never visible strings. If the application
records which commands are used, record \api{ActionKey} ids. If it records
which help topics open, record \api{HelpTopic} ids. If it records validation
friction, record problem codes and field keys. Visible text changes because
the product improves. Analytics should not split a time series because
\api{"Save customer"} was copyedited to \api{"Save changes"}.

Migration from older Swing code should start with operations that already
appear in several places. Save, Refresh, Delete, Export, Search, Apply,
Cancel, and Help are good candidates. Introduce action keys and descriptors.
Create commands. Adapt existing buttons and menu items. Move enablement into
the command. Then migrate shortcuts. Only after the operation path is stable
should you clean up resource keys and help topics. This sequence reduces risk:
the user sees the same screen while the internal identity improves.

Do not migrate only labels. Replacing string literals with resource keys is
useful, but it is not the whole chapter. A screen can have resources and still
have duplicated listeners, untyped help topics, tooltip races, and tests that
search for visible text. The identity review must cover operations, fields,
components, resources, help, and tests together. Otherwise the screen gains a
resource bundle but remains structurally fragile.

Generated metadata needs review because generated mistakes are multiplied.
Check generated action ids for consistency. Check resource ids for predictable
families. Check help topics for durable names. Check descriptor lists for
expected order. Check selectable flags for toggles. Check processor diagnostics
for duplicates. A generated companion should make the screen easier to audit,
not merely produce more files. If generated names are hard to predict, the
book's examples and the project's tests will both become harder to read.

Handwritten metadata should not be treated as second-class. Some screens will
never be generated because they are too specialized, too dynamic, or too small.
They should still use the same typed key families. A handcrafted GIS toolbar,
paint editor, or diagnostic console can declare action keys, command
descriptors, help topics, and component keys. Then the rest of Corusco's Swing
infrastructure can test and bind it the same way as generated screens.

Finally, teach the vocabulary to the team. During review, say ``this is an
action key'', ``this is a resource key'', ``this is a help topic'', ``this is
a component key'', and ``this is visible text.'' The words prevent ambiguity.
Without them, developers say ``id'' and mean five different things. Once the
vocabulary is shared, many design decisions become mechanical: stable identity
goes in keys, presentation goes in resources and components, behavior goes in
commands and bindings, help goes through topics, tests use typed handles.

\section{Exercises and checklist}

\begin{exercisebox}
\begin{enumerate}
\item Find a screen with duplicated button and menu listeners. Replace them
with one command and adapt it to both surfaces.
\item Add an \api{ActionDescriptor} with text, tooltip, mnemonic, and
accelerator resource metadata for one operation.
\item Convert a check-box menu item and toolbar toggle button to share one
selectable command or model value.
\item Add a \api{CommandSet} to a presenter and update tests to look up
commands by \api{ActionKey}.
\item Replace one button-text lookup in a test with a \api{ComponentKey} or
\api{ActionKey} lookup.
\item Add required resource tests for generated command text, field labels, and
table headers.
\item Add a help topic to a field or table column and test that F1 dispatches a
structured \api{HelpRequest}.
\item Replace competing calls to \api{setToolTipText} with one composed
tooltip policy.
\item Review generated action metadata and decide which parts are descriptors,
which parts are runtime command state, and which parts are Swing bindings.
\item Localize or copyedit one visible string and verify that tests and
persistence still use stable ids rather than visible text.
\end{enumerate}
\end{exercisebox}

\begin{center}
\textbf{Identity checklist.}
\end{center}

\begin{itemize}
\item Repeated operations have \api{ActionKey} values and commands.
\item Command descriptors use resource keys, not hard-coded visible text.
\item Buttons, menus, toolbars, popup items, and shortcuts share command state.
\item Toggle state has one owner.
\item Disabled reasons are modeled when users need explanation.
\item Field, component, resource, and help identities are separate typed keys.
\item Required resources are verified through typed keys.
\item Help topics are stable ids, not tooltip text.
\item Tooltips compose problems, disabled reasons, static help, and help hints.
\item F1/help behavior is installed and cleaned up through a lifecycle scope.
\item Tests use \api{ActionKey}, \api{FieldKey}, and \api{ComponentKey} where possible.
\item Localized strings are presentation, not persistence or test identity.
\item Generated companions expose stable facts without hiding screen design.
\item Swing adapters and behaviors are closed with the view that installed them.
\end{itemize}

The purpose of all this naming is not ceremony. It is to keep a Swing
application from forgetting what it already knows. A command is one operation
even when it appears in five places. A resource is one typed value even when it
is resolved through several locales. A help topic is one durable help target
even when the help viewer changes. A field is one semantic fact even when its
component moves. Once identity is stable, presentation can improve freely.
